\section{Results}
\label{sec:results}

\subsection{The mechanism is far from inert}
\label{sec:res:magnitude}

Before asking whether the centrality weighting is informative we establish that
it is active, since a control on an inert mechanism is uninformative
(Section~\ref{sec:leverage}). Fig.~\ref{fig:magnitude} reports the normalised
displacement $\rho$ of \eqref{eq:relshift} for the graph configurations used in
the original report and in its reference implementation.

\begin{figure}[t]
\centering
\includegraphics{fig_mechanism_magnitude}
\caption{Displacement induced by centrality weighting, in units of the
population's own spread. The mechanism moves each update by a quarter to a
third of the population spread; it is not a small correction.}
\label{fig:magnitude}
\end{figure}

Across the published and implemented settings $\rho$ ranges from $0.25$ to
$0.31$: the centrality-weighted aggregate differs from the plain neighbourhood
mean by roughly a third of the population's spread. The coefficient of
variation of betweenness over nodes is $0.65$--$0.89$, so the weights are
strongly heterogeneous rather than near-uniform.

Two predictions of Section~\ref{sec:leverage} are confirmed directly. On the
ring lattice ($p=0$), which is vertex-transitive, we measure
$\rho = 3.1\times10^{-17}$: the mechanism is inert by construction, exactly as
$\chi_i \equiv 0$ requires. And because betweenness is exactly zero at many
nodes of a sparse small-world graph, the guard $\varepsilon$ drives the
within-neighbourhood weight ratio to $O(10^{6})$--$O(10^{7})$: the operator
behaves as hard masking rather than soft weighting.

\subsection{The assignment carries no information}
\label{sec:res:placebo}

Table~\ref{tab:placebo} gives the comparison between
the published method and its permutation control.

\begin{table}[t]
\caption{\textsc{Social} against its permutation control. Paired Wilcoxon per
function with Holm correction; $\bar{\delta}$ is the mean effect over the
suite and the interval is its $90\%$ bootstrap CI. No comparison yields a
single significant difference.}
\label{tab:placebo}
\centering
\small
\setlength{\tabcolsep}{4pt}
\begin{tabular}{@{}lccccl@{}}
\toprule
Suite & $D$ & Budget & $n_f$ & W/L & $\bar{\delta}$ [90\% CI] \\
\midrule
Classic & 30 & $2.0\!\times\!10^{4}$  & 10 & 0/0 & $+0.048$ \\
Classic & 30 & $1.05\!\times\!10^{5}$ & 10 & 0/0 & $+0.078$ \\
CEC2017 & 10 & $1.0\!\times\!10^{5}$  & 28 & 0/0 & $-0.004$ $[-0.034,+0.026]$ \\
CEC2017 & 30 & $3.0\!\times\!10^{5}$  & 28 & 0/0 & $-0.010$ $[-0.038,+0.020]$ \\
CEC2017 & 50 & $5.0\!\times\!10^{5}$  & 28 & 0/0 & $-0.032$ $[-0.066,+0.002]$ \\
\bottomrule
\end{tabular}
\end{table}


On $104$ function comparisons spanning two suites, three dimensions, five
budgets and a twenty-fivefold range of evaluation counts, permuting betweenness
values across nodes changes nothing that survives correction. The largest
effect anywhere is $|\delta|=0.26$; suite means are within $\pm0.08$ of zero. A
mechanism that displaces every update by a third of the population spread can
have its node-to-score assignment randomised without measurable consequence.

The dimensional invariance is a prediction of the framework rather than an
additional observation. By \eqref{eq:leverage} the structural leverage
$\chi_i$ depends only on the graph and the score vector, not on $D$, so the
mechanism's strength---and hence what its randomisation can destroy---should
not change with dimension. Across the three protocol dimensions the suite-level
effect is $-0.004$, $-0.010$ and $-0.032$, all equivalent to zero within
$|\delta|<0.10$ and the first two within $|\delta|<0.05$.

\subsection{Equivalence, and the limits of the present design}
\label{sec:res:tost}

A non-significant difference is not equivalence, so we test the latter
directly (Section~\ref{sec:stats}). Here the answer depends on the level of
aggregation, and we report both.

At the level of the suite, equivalence holds comfortably, and tightens with
sample size. At $D=30$ with the $51$ runs the CEC2017 protocol prescribes, the
mean effect over the $28$ functions is $\bar{\delta}=-0.010$ with a $90\%$
interval of $[-0.038,+0.020]$: contained not only in the negligible band but in
the far stricter $|\delta|<0.05$. A one-sample Wilcoxon test on the
per-function effects gives $p=0.34$. At $D=10$ with $30$ runs the mean effect
is $-0.004$, interval $[-0.034,+0.026]$, likewise inside $|\delta|<0.05$.

Per function the picture is different, and the reason is worth separating into
two parts. Raising the runs from $30$ to $51$ narrows the median $90\%$
half-width from $0.181$ to $0.139$---exactly the $\sqrt{30/51}$ the sampling
theory predicts---and raises the count of formally equivalent functions from
$1$ to $6$ of $28$. It does not raise it further, because equivalence at margin
$m$ requires $|\hat{\delta}|+\text{half-width}<m$, not merely
$\text{half-width}<m$: with a median $|\hat{\delta}|$ of $0.063$, reaching
$|\delta|<0.147$ on the median function would take roughly $100$ runs. And no
sample size suffices for all of them: on $4$ of $28$ functions the estimated
effect already exceeds $0.147$, so a genuine, if small, difference is present
there. At the small-effect margin $|\delta|<0.33$, $26$ of $28$ are equivalent.

Per-function intervals are given in the supplementary material. Most sit
inside the band but cross its edge, which is the visual form of this
distinction: failing to detect a difference is not the same as establishing
its absence.


We state the conclusion at the strength the evidence supports, which is not
that the two arms are identical. Across the suite the assignment of centrality
scores to graph positions has an effect equivalent to zero within a margin four
times stricter than the negligibility threshold. Per function, effects are
bounded within a small margin, a handful are near the negligibility boundary,
and none survives correction as a difference.

\subsection{Heterogeneity does what the assignment does not}
\label{sec:res:p1p2}

Section~\ref{sec:claim} separated two properties that ablation confounds. The
data separate them too. Comparing the permutation control against uniform
mixing on the same graph isolates (P1), heterogeneity, with the assignment
already destroyed: \permcontrol{} beats \textsc{Unif}$_4$ on $3$ of $28$
CEC2017 functions ($F_{12}$, $F_{22}$, $F_{23}$; $\delta=-0.37,-0.53,-0.74$;
Holm-corrected $p\le0.013$) and loses on none. Comparing \textsc{Social}
against \permcontrol{} isolates (P2), the assignment: $0$ of $28$.

Unequal weights therefore help slightly; \emph{which} neighbour receives the
large weight does not matter at all. An ablation reports the sum of the two and
attributes it to structure. This is the confound of Section~\ref{sec:claim},
measured.

The practical consequence is immediate. Betweenness is roughly $30\%$ of
per-iteration runtime and can be replaced by any heterogeneous weighting---a
fixed random draw suffices---at no cost in solution quality.

\subsection{Calibrating the instrument}
\label{sec:res:calibration}

A null result is worth only as much as the test's ability to return a non-null
one. We therefore ran the same controls on two further mechanisms of the same
optimizer, on the same suite, at the same sample size and with the same
statistic. Fig.~\ref{fig:calibration} places the three side by side.

\emph{Fitness influence} (Section~\ref{sec:positive}) is the positive control:
the permuted score is a direct function of objective values the run has already
paid for. Permuting it does cost performance, but less uniformly than we
expected. At suite level the effect is $\bar{\delta}=-0.051$ with $90\%$
interval $[-0.117,+0.008]$, which includes zero; $17$ of $28$ functions favour
the unpermuted arm. The effect is concentrated by problem class: $-0.150$ on
the nine composition functions and $-0.152$ on the two unimodal ones, against
$+0.012$ on hybrids and $+0.016$ on simple multimodals. On $F_{23}$ it reaches
$\delta=-0.80$ with Holm-corrected $p=0.003$.

That single function settles the power question for the permutation control:
the same procedure, applied to a different score on the same optimizer, detects
an effect of $0.80$. It is not blind. Applied to the centrality assignment it
returns $-0.010$ with an interval four times inside the negligibility margin.

\emph{Adaptive selection} is the policy control applied to \textsc{Ads}, which
chooses neighbourhood degree online by sliding-window UCB on observed
improvement. Here the instrument returns a clear signal: $\bar{\delta}=-0.101$
with interval $[-0.149,-0.061]$ excluding zero, a one-sample Wilcoxon
$p=1.5\times10^{-5}$, and $24$ of $28$ functions favouring informed selection
over uniform switching.

Two points follow, and the second corrects an earlier reading of our own
results. First, the ordering across the three contrasts---null, intermediate,
clear---is exactly what a working instrument should produce, and it is produced
without any change of suite, sample size or statistic. Second, per-function
Holm correction across $28$ comparisons is conservative enough to hide a real
suite-level effect: it credits \textsc{Ads} with wins on only $3$ functions,
which reads as a weak effect, while the suite-level estimate is unambiguous.
Where a mechanism's benefit is diffuse rather than concentrated, the
per-function view understates it, and both should be reported.

\begin{figure}[t]
\centering
\includegraphics{fig_calibration}
\caption{The same controls on three mechanisms of the same optimizer, on
CEC2017 at $D=30$. Points are per-function effects, intervals are $90\%$
bootstrap CIs on the suite mean. The instrument separates a null mechanism from
an intermediate and a clear one without changing suite, sample size or
statistic.}
\label{fig:calibration}
\end{figure}

Mean ranks over the six arms at $D=30$ place the two sparse arms first
(\textsc{Social} $3.121$, \permcontrol{} $3.170$), then uniform mixing
($3.335$), with the density-varying arms last ($3.676$--$3.977$). That
\textsc{Social} and \permcontrol{} are adjacent, separated by $0.05$ rank units
with overlapping standard errors, is the finding of
Section~\ref{sec:res:placebo} in another form. Per-function medians and the
rank plot are given in the supplementary material.


\subsection{What density does}
\label{sec:res:density}

A sweep over $K\in\{4,8,16,32,59\}$ across $49$ problem configurations at their
native dimensions, $7{,}350$ runs, finds neighbourhood size to matter in $30$
of $49$ (Friedman $p<0.05$), with rank spreads up to $3.9$ of a possible $5$
and $p$ as small as $2\times10^{-23}$. Where it matters the effect is large;
in the remaining $19$ configurations it is undetectable.

The optimum is problem- and dimension-dependent. Ruggedness, measured as the
lag-one autocorrelation of fitness along a random walk and therefore computable
online, correlates with the preference for dense neighbourhoods at
$\rho_{s}=-0.61$ ($p=3.7\times10^{-4}$) over the configurations where density
matters (supplementary material).

That correlation is carried largely by the three most rugged
functions, and it vanishes at $D=10$ ($\rho_{s}=-0.15$, $p=0.63$) while holding
at $D=30$ ($-0.73$) and $D=50$ ($-0.63$). \textsc{Penalized} and
\textsc{Penalized2} are clear counterexamples, smooth by the proxy yet
preferring the densest neighbourhood at every dimension.

Density is thus consequential, not predictable in advance from this
descriptor, and not constant across dimension---an argument for adapting it
rather than fixing it by a topology prior, and simultaneously an argument
against predicting it from landscape features.
